\title{NeighborCache: Caching the SSD on the SSD for Graph-Based Vector Search}

\begin{abstract}
Disk-based graph indexes such as DiskANN make billion-scale vector search affordable, but each query must fetch tens to hundreds of pages from the SSD, and these reads dominate its latency.
A read pays off only when the fetched page also holds the nodes the query will visit next, which in a graph traversal are the neighbors of the node just read.
Proximity graphs, however, give every node dozens of neighbors: no arrangement of a single copy can keep a node beside all of them, and replicating every node's neighborhood multiplies the index size.
Existing systems therefore settle for one copy, caching nodes near the entry point or reordering the graph on disk, and they make these choices from the index before serving.
We observe that queries need only a small part of this locality: under skewed workloads, searches concentrate on a few hot nodes and repeatedly follow the same few edges out of them.
We present \textsc{NeighborCache}, which caches exactly this part, and caches it on the SSD itself.
\textsc{NeighborCache} tracks the hotness of nodes and edges during search, keeps the hottest nodes in memory, and copies each hot node together with the neighbors reached through its hot edges into a dedicated SSD page, so that one read replaces several reads to the original index.
The gain comes from fewer reads rather than faster ones, so the cache can live on the same SSD it accelerates.
Cache pages are capped at 20\% of the index and leave the index and its search algorithm untouched; as hotspots shift, stale hotness decays and pages are replaced online without rebuilding the index.
A comprehensive experimental study demonstrates that \textsc{NeighborCache} improves search throughput by up to 84\% over DiskANN while maintaining the same recall and memory budget, and reduces disk reads per query by up to 55\%; the gain already appears on the dataset's nearly uniform query set and grows as the workload becomes more skewed or repetitive.
\end{abstract}
